\section{六次Gram曲線からの復元}
\label{sec:prym-torelli-sextic-v26}

署名 $(3,3)$ の成分 $Z_1$ では、二つの固有空間の順序を保つ。
すなわち $P$ に用いる三次元固有空間を指定し、
関係基底変換、二つのブロック内の基底変換、外側の座標変換による有理商を
$\mathfrak Z_1^+$ と書く。
その次元は命題\ref{prop:prym-baseline-v26}と同じく七である。
まずこの順序を保持して逆構成を記述する。
第\ref{subsec:unordered-blocks-v27}項では、順序を忘れた商も同じ関数体を持つことを示す。

\begin{theorem}[平面Prymデータによる復元]
\label{thm:plane-prym-torelli-v26}
$\mathfrak Z_1^+$ の非空開部分では、平面Prymデータ
\[
 (C,\eta,H),\qquad
 C=V(\det(PAP^t)),\quad H=\mathcal O_C(1),\quad
 \eta=\operatorname{coker}(PAP^t)(-2)
\]
が、指定した順序を保つ同型を除いて対象を一意に復元する。
従ってこの写像は像への双有理写像であり、その像の次元は七である。
平面直線束の標識は、第\ref{subsec:abstract-prym-torelli-v27}項で除く。
指定反対合そのものの標識は保持する。
\end{theorem}

\begin{proof}
まず、既存の $Z_1$ の非空開部分で次の二つの階数条件が成立することを確認する。
$P(x)=\sum_{i=0}^2x_iP_i$ の九つの係数行ベクトルを並べた
$9\times7$ 行列を $\widehat P$ とする。
また
\[
 K_{P,A}=\{b(x)\in k^7\otimes H^0(\mathcal O(1)):
                              P(x)A b(x)^t=0\}
\]
と置く。
第\ref{sec:gram-prym-v25}節と同じ素数 $101$、seed $17029$ のHensel点で、
以下の小行列式はいずれも $101$ 進単元である。
\[
\begin{array}{c|c|c}
 \text{行列}&\text{階数}&\text{指定小行列式の }101^3\text{ における剰余}\\\hline
 \widehat P\ (9\times7)&7&755782\\
 \widehat P A\widehat P^t\text{ の指定主 }7\times7\text{ 部分行列}&7&389051\\
 PAb^t=0\text{ の係数行列 }(18\times21)&18&990506\\
 B\text{ の三行の }K_{P,A}\text{ 内での座標行列}&3&60902
\end{array}
\]
従って $\operatorname{rank}\widehat P=7$、$\dim K_{P,A}=3$、
そして $B$ の三行は $K_{P,A}$ 全体の基底をなす。
これは $101$ 進点だけの性質ではなく、同じ既約成分上の非空開条件である。
この開部分を、二十一点の被約性と $C$ の滑らかさの既存の非空開部分と交わらせる。

次に $Q=PAP^t$ から $(P,A)$ を復元する。
$Q_{uv}$ の $x_i x_j$ の係数を $q_{uv,ij}$ とし、
九行九列の行列 $G(Q)$ を
\begin{equation}\label{eq:sextic-coefficient-gram-v26}
 G(Q)_{(u,i),(v,j)}=
 \begin{cases}
  q_{uv,ii},&i=j,\\
  q_{uv,ij}/2,&i\ne j
 \end{cases}
 \qquad (0\le u,v,i,j\le2)
\end{equation}
で定める。添字を入れ替える場合には $Q_{vu}=Q_{uv}$ とする。
同符号ブロックの全incidenceは
\[
 (P_i)_u A(P_j)_v^t=(P_j)_u A(P_i)_v^t
\]
を含むため、式\eqref{eq:sextic-coefficient-gram-v26}は
\[
 G(Q)=\widehat P A\widehat P^t
\]
を正確に復元する。
特にその階数は七である。

可逆な主部分行列 $G_I=G(Q)_{I,I}$ を選び、
\begin{equation}\label{eq:sextic-rational-factorization-v26}
 A_0=G_I,\qquad
 \widehat P_0=G(Q)_{*,I}G_I^{-1}
\end{equation}
と置く。
階数七のSchur補行列が零になることから
$G(Q)=\widehat P_0A_0\widehat P_0^t$ である。
これは $Q$ の係数の有理式による復元である。
他の非退化な階数七の分解
$G(Q)=\widehat P_1A_1\widehat P_1^t$ があれば、
$\operatorname{im}\widehat P_1=\operatorname{im}G(Q)
=\operatorname{im}\widehat P_0$ なので
\[
 \widehat P_1=\widehat P_0g,\qquad
 A_1=g^{-1}A_0g^{-t}
 \quad(g\in\operatorname{GL}_7)
\]
となる。従って $(P,A)$ は関係基底変換だけの差を除いて一意である。

復元された $(P,A)$ から十八行二十一列の線形系
$PAb^t=0$ を作る。その核は三次元であり、$B$ の三行はこの核全体を張る。
従って $B$ も $\operatorname{GL}_3$ による行変換だけの差を除いて一意である。
ここで $Q$ は元の成分の像に属するものを扱っているので、
復元した核の基底が満たす $B$--$B$ ブロックのincidenceは元の対象から従う。
任意の階数七の $G(Q)$ についてそれらの方程式が自動的に成立するとは主張しない。

最後に平面Prymデータから $Q$ を復元する。
このデータは $\eta(2H)$ とその最小次数付き自由分解
\[
 0\longrightarrow\mathcal O_{\mathbb P^2}(-2)^3
 \xrightarrow{Q}\mathcal O_{\mathbb P^2}^{3}
 \longrightarrow i_*\eta(2H)\longrightarrow0
\]
を同型を除いて決める。
対称な最小分解の一意性により、二つの対称表示は合同変換で結ばれる
\cite[\S2.3]{BeauvilleDeterminantal}。
従って $(C,\eta,H)$ は $Q$ を正ブロックの $\operatorname{GL}_3$ と
平面の $\operatorname{PGL}_3$ の作用だけの差を除いて決める。
前二段階の復元と合わせて、$\mathfrak Z_1^+$ の一般点は一意に決まる。
標数零の既約な粗有理商の間で、一般の幾何学的ファイバーが一点なので、
像への写像は双有理である。像の次元は従って七である。

\end{proof}

\paragraph{Gram像の次元の独立検算.}
同じHensel点の正規化チャートは百五十三変数、七十二方程式を持ち、接次元は八十一である。
三十六個のGram係数の微分をこの接空間へ制限すると、指定二十四次小行列式が単元となる。
その $101^3$ における剰余は $822512$ であり、従って像の次元は少なくとも二十四である。
一方、正規化前の全pairingチャートは八十二次元であり、
関係基底の $\operatorname{GL}_7$ と負ブロックの $\operatorname{GL}_3$ は
$Q$ を変えない五十八次元の自由な作用を持つ。
自由性は $\operatorname{rank}\widehat P=7$ と $B$ の行独立性から従う。
従ってGram像の次元は高々 $82-49-9=24$ であり、正確に二十四である。
正ブロックと座標の二つの $\operatorname{GL}_3$ の作用は、
互いに逆のスカラーによる一次元核を持ち、有効軌道の次元は十七である。
従って商の次元は $24-17=7$ となり、上の復元による次元計算と一致する。

\paragraph{再現入口.}
\nolinkurl{gram_dominance_v26.py} を引数 \texttt{--signature 3} で実行する。
対応する \nolinkurl{gram_dominance_v26_p3_certificate.json} は、
階数証明書、線形核の基底、Gram分解の有理式の照合、
二十四次微分小行列式を含む。
全ての入力精度は $101^{32}$ であり、ここで使う小行列式は全て単元である。
従って階数と復元の成立は既存の同じHensel点で検証されている。
